% =============================================================================
%  CAPÍTULO 5 — CONCLUSÃO
%  TCC: Dollar Bars vs. Time Bars — Modelos GARCH para BTC/USDT
% =============================================================================

\chapter{Conclusão}
\label{cap:conclusao}

% -----------------------------------------------------------------------------
\section{Retomada dos Objetivos Específicos}
\label{sec:retomada_objetivos}

Esta seção retoma os cinco objetivos específicos enunciados na seção 1.1.2 e avalia
sua consecução à luz dos resultados obtidos.

O primeiro objetivo — construir séries de retornos a partir de time bars e dollar bars
e comparar suas propriedades estatísticas — foi plenamente atingido. As dollar bars
produziram retornos com curtose de Fisher de 3,09, contra 75,37 das time bars, uma
redução de cerca de 25 vezes. A assimetria das dollar bars foi de 0,05, próxima de zero,
enquanto as time bars registraram $-1{,}01$. Esses resultados confirmam quantitativamente
as previsões de \citeonline{lopez2018advances} sobre o efeito do mecanismo de amostragem nas
propriedades distribucionais dos retornos.

O segundo objetivo — estimar a variância realizada por meio do estimador TSRV — também
foi atingido. O estimador \textit{Two-Scale Realized Variance} foi calculado sobre dados
\textit{tick-by-tick} com três frequências de subamostragem (30 segundos, 1 minuto e
2 minutos), e a média dos três estimadores foi utilizada como benchmark de volatilidade
diária. O procedimento de múltiplas escalas reduz o viés introduzido pelo ruído de
microestrutura, produzindo um estimador consistente da variância quadrática integrada
\cite{zhang2005}.

O terceiro objetivo — ajustar modelos ARMA-GARCH e selecionar ordens ótimas via critérios
de informação — foi atingido. O CPCV com 252 folds selecionou o modelo ARMA$(0,0)$-GARCH$(1,2)$
para as dollar bars, com frequência de 92,5\% dos folds (233/252), e o modelo
ARMA$(0,0)$-GARCH$(1,1)$ para as time bars, com frequência de 100\% dos folds (252/252).
Em ambos os casos, a seleção foi estável e robusta ao longo dos folds.

O quarto objetivo — avaliar e comparar o desempenho preditivo com MSE, MAE, QLIKE e
a regressão de Mincer-Zarnowitz — foi atingido. As dollar bars dominaram em todas as
quatro dimensões de avaliação. O coeficiente de determinação MZ foi de $R^2 = 0{,}655$
para as dollar bars, contra $R^2 = 0{,}493$ para as time bars, diferença de 16,2 pontos
percentuais.

O quinto objetivo — comparar o desempenho preditivo tendo o estimador TSRV como
referência — foi igualmente atingido. O TSRV serviu como benchmark de volatilidade diária;
ambos os modelos GARCH apresentaram subestimação sistemática leve, evidenciada pelos
interceptos negativos da regressão MZ ($\hat{\alpha}_{\text{DB}} = -1{,}08 \times 10^{-4}$
e $\hat{\alpha}_{\text{TB}} = -2{,}54 \times 10^{-4}$). Esse padrão é característico de
modelos GARCH paramétricos estimados com inovações gaussianas quando confrontados com
estimadores não-paramétricos de alta frequência, que capturam componentes de variação
intrajornada não presentes no modelo.

% -----------------------------------------------------------------------------
\section{Resposta à Pergunta Central}
\label{sec:resposta_central}

A pergunta de pesquisa que motivou este trabalho indaga: \textit{em que medida a escolha
do método de amostragem afeta a qualidade das estimativas de volatilidade obtidas por
modelos ARMA-GARCH?}

Os resultados indicam que a escolha afeta de forma relevante tanto as propriedades
estatísticas das séries de retornos quanto o desempenho preditivo fora da amostra dos
modelos. No plano distribucional, a diferença é dramática: curtose 25 vezes menor e
assimetria quase nula nas dollar bars, frente a caudas extremamente pesadas e assimetria
negativa pronunciada nas time bars. No plano preditivo, a diferença é economicamente
significativa: redução de 34\% no MSE, de 24\% no MAE e de 16 pontos percentuais no
$R^2$ da regressão de Mincer-Zarnowitz, além de ICs bootstrap para o QLIKE que não
se sobrepõem a 95\% de confiança.

A significância estatística formal, avaliada pelo teste de Diebold-Mariano, não atinge
o limiar convencional de 5\% (estatística DM = $-1{,}9578$, p-valor = 0,0503), mas
o argumento baseado exclusivamente nessa fronteira seria epistemicamente incompleto.
O poder do teste é limitado pelo número de dias de sobreposição entre os conjuntos de
previsão (147 dias), e o p-valor de 0,0503 constitui evidência fraca contra a hipótese
nula, não evidência a seu favor. A combinação de sinal inequívoco na estatística DM,
dominância em todas as métricas pontuais e não sobreposição dos ICs bootstrap para o QLIKE
configura evidência consistente — ainda que não conclusiva ao nível convencional — em
favor das dollar bars.

É importante destacar que este resultado não é apenas confirmatório de \citeonline{lopez2018advances}.
O autor argumenta, no plano teórico e com ilustrações empíricas, que barras de informação
produzem séries com propriedades distribucional superiores, mas não testa diretamente
modelos de previsão de volatilidade nem utiliza métricas específicas para comparação de
previsões. O presente trabalho estende essa análise para o contexto de modelos GARCH —
a classe de modelos mais amplamente empregada para volatilidade condicional — e quantifica
a diferença de desempenho com funções de perda robustas da classe de \citeonline{patton2011},
fornecendo evidência empírica nova sobre a utilidade prática das dollar bars em um contexto
de previsão.

% -----------------------------------------------------------------------------
\section{Contribuições do Trabalho}
\label{sec:contribuicoes}

Este trabalho apresenta três contribuições concretas à literatura.

A primeira é de natureza metodológica. A utilização do \textit{Combinatorial Purged
Cross-Validation} para seleção de hiperparâmetros em modelos GARCH — procedimento que
elimina vazamento de informação temporal por meio de purga e embargo, e que explora
todas as $\binom{N}{K}$ combinações de folds em vez de um único caminho de validação —
representa um avanço em relação ao walk-forward simples usualmente empregado na literatura
brasileira de modelos de volatilidade. A implementação computacional, que realizou 72.576
ajustes MLE em aproximadamente 17 minutos por meio de um backend em Rust, é documentada
de forma reproduzível e pode ser reutilizada para outros ativos e especificações.

A segunda contribuição é empírica. O presente trabalho oferece a primeira comparação
sistemática, ao menos no contexto da literatura acadêmica brasileira, entre dollar bars
e time bars como insumo para modelos GARCH aplicados a criptomoedas, com avaliação
out-of-sample via estimador TSRV e métricas robustas de Patton (2011). A extensão da
análise de \citeonline{lopez2018advances} para o problema específico de previsão de volatilidade
condicional preenche uma lacuna identificada na revisão de literatura.

A terceira contribuição é técnica. O \textit{pipeline} computacional desenvolvido em
Python com backend em Rust demonstra que é possível realizar seleção de hiperparâmetros
por CPCV para modelos GARCH em escala prática — 36.288 ajustes MLE por série em menos
de nove minutos por pipeline —, tornando o procedimento viável sem infraestrutura
computacional especializada.

% -----------------------------------------------------------------------------
\section{Limitações}
\label{sec:limitacoes_conclusao}

As principais limitações do trabalho foram discutidas em detalhe na seção 4.4.
Resumidamente: os resultados referem-se a um único ativo de criptomoeda (BTC/USDT na
Binance), sobre um período de dois anos (2024–2025), com limiar $\bar{D}$ fixo calibrado
para equiparar o número médio de barras por dia entre as duas séries, e com inovações
gaussianas em ambos os modelos GARCH. A generalização das conclusões para outros ativos,
períodos ou especificações requer estudos adicionais.

% -----------------------------------------------------------------------------
\section{Agenda de Pesquisa Futura}
\label{sec:agenda_futura}

Os resultados obtidos abrem um conjunto de direções para investigação futura.

A mais imediata é a replicação do estudo em outras classes de ativos. Ações brasileiras
listadas na B3, o câmbio BRL/USD e commodities negociadas em mercados regulados apresentam
microestruturas distintas das criptomoedas — com leilões de abertura e fechamento,
circuit breakers e menor fragmentação de liquidez — o que permitiria investigar se a
superioridade das dollar bars é específica ao mercado de criptomoedas ou se é uma
propriedade mais geral do método de amostragem.

Uma segunda extensão natural é a re-estimação dos modelos com distribuições de cauda pesada
para as inovações, particularmente a t-Student e a Distribuição de Erro Generalizado (GED).
Dado o excesso de curtose documentado, especialmente nas time bars, a adoção de inovações
não-gaussianas poderia melhorar o ajuste e, potencialmente, alterar a magnitude das
diferenças de desempenho entre os tipos de barra.

\citeonline{lopez2018advances} propõe outros tipos de barras de informação além das dollar bars,
entre elas as volume bars e as tick bars, bem como as imbalance bars, que condicionam o
tamanho da barra ao desequilíbrio entre ordens de compra e venda. Uma comparação mais
abrangente entre esses tipos de barra forneceria um panorama mais completo sobre qual
estrutura de amostragem é mais informativa para modelos de volatilidade.

No plano dos modelos de volatilidade, caberia investigar se a superioridade das dollar bars
se mantém com especificações que capturam assimetria e efeitos de alavancagem, como o
EGARCH \cite{Nelson1991} e o GJR-GARCH \cite{GlostenJagannathanRunkle1993}, ou com
modelos de volatilidade estocástica, que tratam a variância latente como uma variável
de estado não-observada. Esses modelos podem ser particularmente relevantes para as
dollar bars, cuja dinâmica de volatilidade demandou um componente GARCH$(1,2)$ em vez
do canônico GARCH$(1,1)$.

Por fim, o presente trabalho avaliou previsões \textit{one-step-ahead} da variância
diária. Previsões multi-step — de 5, 10 e 22 dias úteis, horizontes relevantes para
gestão de risco e cálculo de capital regulatório — podem exibir padrões distintos de
desempenho relativo entre os tipos de barra, em função da estrutura de persistência da
volatilidade capturada por cada modelo. Essa extensão representaria uma contribuição
diretamente aplicável a contextos práticos de gestão de risco quantitativo.
